\chapter{Research Methodology}
\label{ch:methodology}

\begin{chaptersummary}
This chapter sets out the design method that decides where a model belongs in the agent and what it may
decide, defines the ACRS closed loop and the episode state shared by the four objectives, and describes
three evaluation layers: sessionised ToolBench with simulated agents for controlled, comparable
measurement; KCC-Sessions, built from real farmer queries, for the public problem; and AppWorld with
real open-weight agents for external validity. It then gives the leakage audit, the protocol, the
metrics and the statistical analysis, including the sample sizes each claim needs.
\end{chaptersummary}

\section{Design Method}
\label{sec:method-design}

Every component of ACRS was designed by the same four-phase method. \emph{Understand} traces a symptom,
such as a low score or a high cost, back to its root cause. \emph{Define} states the target outcome
independently of the metric one is asked to move: for ACRS, \emph{every auto-executed action is
grounded, affordable and within a certified error bound, and everything else reaches a named person with
its evidence attached}. \emph{Decompose} asks of every step, in a fixed order, whether it involves
judgment at all (if not, it is plain code), whether its steps are knowable in advance (a workflow) or
whether a model must choose what happens next (an agent, which then needs a termination rule); and then,
separately, what the model may decide once it is there. \emph{Sequence and validate} stages deployment
so that each stage produces the evidence for the next: shadow mode, then live with approval, then a
separately reviewed path for any change of scope. The resulting stage registry, with each stage's
pattern, authority and failure detector, is given in Chapter~\ref{ch:governance}.

\section{The ACRS Closed Loop}
\label{sec:method-loop}

ACRS processes each step through four stages in a fixed order and writes the result back into a state
that the next step reads (Figure~\ref{fig:acrs-loop}):
\begin{equation}
q_t,\ M_{t-1} \;\xrightarrow{\ \text{SATR}\ }\; \mathcal{T}_t
\;\xrightarrow{\ \text{APRR}\ }\; a_t
\;\xrightarrow{\ \text{MNCD}\ }\; \mathcal{C}_t
\;\xrightarrow{\ \text{FCNP}\ }\; (\hat{y}_t,\ M_t),
\label{eq:acrs-loop}
\end{equation}
where $q_t$ is the request or sub-goal at step $t$, $\mathcal{T}_t$ the ranked tool shortlist, $a_t$ the
agent or model tier chosen, $\mathcal{C}_t$ the consensus certificate, and $\hat{y}_t$ the answer or
action (Figure~\ref{fig:turn-loop}). A governance stage G1 converts $\mathcal{C}_t$ into a disposition
(execute, review or block) by a deny-by-default policy check. Integration is not a fifth algorithm: it
is the contract that every stage reads from and writes to the same state.

\begin{figure}[htbp]
\centering
\includegraphics[width=\textwidth]{figures/02_system_architecture}
\caption{ACRS system architecture: the four objectives share the episode state $M_t$ and draw tools from
the active environment (ToolBench, Indian OGD services or AppWorld applications)}
\label{fig:acrs-loop}
\end{figure}

\begin{figure}[htbp]
\centering
\includegraphics[width=0.8\textwidth]{figures/03_turn_loop}
\caption{One step of the ACRS closed loop, from the request to the write-back of $M_t$}
\label{fig:turn-loop}
\end{figure}

\subsection{Session state}
\label{sec:method-state}

\begin{definition}[Session state]
\label{def:state}
The state after step $t$ is the tuple
\[
M_t = \big(\mathbf{m}_t,\ \mathbf{s}_t,\ \mathbf{f}_t,\ U_t,\ F_t,\ \rho_t,\ K_t\big),
\]
where $\mathbf{m}_t$ is an exponentially decayed term vector of intents (decay $0.72$ per step),
$\mathbf{s}_t$ the term vector of intents already served (decay $0.9$), $\mathbf{f}_t$ the decayed counts
of providers or applications whose APIs were committed (decay $0.8$), $U_t$ and $F_t$ the per-tool usage
and failure counts, $\rho_t \in [0,1]$ the running risk, and $K_t$ the context blocks retained by FCNP.
\end{definition}

Risk rises by $0.18$ on every tool failure and falls by $0.08$ on every success, clipped to $[0,1]$. SATR
reads $\mathbf{m}_t$, $\mathbf{s}_t$, $\mathbf{f}_t$, $U_t$ and $F_t$; APRR reads $\rho_t$ and the remaining
budget; MNCD writes successes, failures and evidence; FCNP writes $K_t$.

\section{Datasets and Evaluation Layers}
\label{sec:method-data}

No single benchmark tests all four objectives with real agents, a public problem and comparability with
prior work, so ACRS is evaluated in three layers (Table~\ref{tab:layers}).

\begin{table}[htbp]
\centering
\caption{The three evaluation layers}
\label{tab:layers}
\footnotesize
\begin{tabular}{@{}p{0.17\textwidth}p{0.22\textwidth}p{0.2\textwidth}p{0.33\textwidth}@{}}
\toprule
Layer & Data & Agents & Purpose \\
\midrule
ToolBench-Sessions & 100 sessions, 174 sub-goals, 10{,}439 APIs & simulated, seeded & controlled ablations, comparability with tool-retrieval literature, adversarial tests \\
KCC-Sessions & real Kisan Call Centre queries; 11 advisory services + ToolBench distractors & simulated, seeded & the public problem; risk classes (informational, dose, entitlement) \\
AppWorld & 9 apps, 457 APIs, 750 tasks with state-based unit tests & real open-weight LLMs (Qwen3, 1.7B--32B) & external validity, cost, collateral damage, leaderboard comparison \\
\bottomrule
\end{tabular}
\end{table}

\subsection{Sessionised ToolBench G1}
The ToolBench G1 retrieval split \citep{qin2024toolllm} has 10{,}439 API endpoints. The 100 test queries
are split into 174 sub-goals, each with its gold endpoint set, and grouped into 100 sessions of one to
four turns. Training trajectories (82{,}599) are used only to fit SATR v1's co-occurrence graph; SATR v2
uses no training data.

\subsection{KCC-Sessions}
Kisan Call Centre transcripts \citep{kccchakshu} hold real farmer queries with expert answers. They become
an ACRS benchmark without inventing labels: the KCC \texttt{QueryType} is mapped once to the advisory
service that answers it (O1 gold); the queries of one block on one day for one crop form a
pseudo-session; \texttt{QueryType} also fixes a risk class (agronomic-dose and entitlement queries are
recommend-only); and the 10{,}439 ToolBench endpoints are added to the corpus as distractors
(Chapter~\ref{ch:governance}).

\subsection{AppWorld}
AppWorld \citep{trivedi2024appworld} provides nine everyday applications operated through 457 APIs and
750 tasks in four splits (\texttt{train}, \texttt{dev}, \texttt{test\_normal} with 168 tasks and
\texttt{test\_challenge} with 417). Each task is graded by unit tests on the final state of the world.
API documentation carries each endpoint's HTTP method, so the side effect of every call is known
exactly. Agents are real open-weight models from the Qwen3 family \citep{yang2025qwen3}, served with vLLM
\citep{kwon2023vllm} on a single free or low-cost GPU; the largest model that fits is used as the large
tier and a 1.7B--4B model as the small tier (Chapter~\ref{ch:appworld}).

\subsection{Indian Open Government Data}
Live execution uses Indian OGD on \texttt{data.gov.in} under the Government Open Data Licence -- India,
including the AGMARKNET daily market-price resource
(\url{9ef84268-d588-465a-a308-a864a43d0070}). Prices are never generated: a failed live call fails loudly
and the turn abstains.

\subsection{Simulated agent pool}
For the ToolBench and KCC layers, agent execution is simulated with a seeded, schema-faithful runtime so
that orchestration mechanisms are measured apart from the quality of any particular LLM. Each agent has a
tier, a cost per call, a latency and a hallucination rate (Table~\ref{tab:agents}). In the adversarial
experiments a hallucinating agent also forges citations and colludes, and poisoned endpoints inject
instructions into tool output.

\begin{table}[htbp]
\centering
\caption{Simulated agent pool (ToolBench and KCC layers)}
\label{tab:agents}
\begin{tabular}{@{}llrrr@{}}
\toprule
Agent & Tier & Cost / call & Latency (ms) & Hallucination rate \\
\midrule
Light lookup & light & 0.0004 & 180 & 0.10 \\
Agriculture specialist & standard & 0.0021 & 420 & 0.04 \\
Geo specialist & standard & 0.0018 & 380 & 0.05 \\
Heavy reasoner & heavy & 0.0135 & 1{,}750 & 0.02 \\
Verifier & verifier & 0.0026 & 310 & 0.01 \\
\bottomrule
\end{tabular}
\end{table}

\section{Leakage Audit}
\label{sec:method-leakage}

Before any comparison, every ToolBench test gold tool set was checked against the training trajectories.
\textbf{94 of the 100 test gold sets appear verbatim in the training data.} A co-occurrence graph fitted
on all trajectories therefore encodes the test answers: with it, SATR v1 reaches NDCG@5 of 0.576; with
the leaking trajectories removed (5{,}476 trajectories, leaving 77{,}123) it falls to 0.485, below the
BM25 baseline of 0.517 (Table~\ref{tab:leakage}). All results use the leak-filtered graph, and the SATR
v2 signals use no training data. On AppWorld the same discipline is enforced by its split rules: prompts
and parameters are tuned on \texttt{train} and \texttt{dev} only, and the test splits are run once per
final configuration without inspecting their tasks.

\begin{table}[htbp]
\centering
\caption{Effect of training-label leakage on SATR v1 (174 sub-goals)}
\label{tab:leakage}
\begin{tabular}{@{}lccc@{}}
\toprule
Configuration & NDCG@5 & NDCG@10 & Recall@10 \\
\midrule
BM25 (stateless) & 0.517 & 0.554 & 0.644 \\
SATR v1, graph fitted on all training data (leaky) & 0.576 & 0.596 & 0.654 \\
SATR v1, leak-filtered graph & 0.485 & 0.522 & 0.592 \\
\bottomrule
\end{tabular}
\end{table}

\section{Evaluation Protocol and Metrics}
\label{sec:method-protocol}

Figure~\ref{fig:method-flow} summarises the evaluation workflow.

\begin{figure}[htbp]
\centering
\includegraphics[width=\textwidth]{figures/08_methodology}
\caption{Evaluation workflow: from the benchmark through the leakage audit to the statistical claims}
\label{fig:method-flow}
\end{figure}

\textbf{Ablation.} Each objective has a switch whose off state is its state-of-the-art analogue: O1
stateless retrieval (BM25 on ToolBench and KCC; static retrieval from the instruction on AppWorld); O2
static routing to the heavy tier; O3 majority vote on ToolBench and KCC and single greedy execution on
AppWorld; O4 a fixed 50\,\% keep ratio or recency truncation. The \emph{SOTA baseline} uses all
analogues; each \emph{+O$k$} configuration swaps in one ACRS mechanism; \emph{Full ACRS} uses all four;
on AppWorld, leave-one-out ablations remove one objective at a time from the full system. On AppWorld the
\emph{baseline} is a ReAct-style code agent that discovers APIs through the documentation app, as in the
benchmark's reference agents.

\textbf{No gold feedback.} SATR commits the system's own picks to the state, never gold tools.

\textbf{Metrics.}
\begin{itemize}[leftmargin=*]
\item Retrieval: NDCG@5 (primary), NDCG@10, Recall@10 and MRR \citep{jarvelin2002cumulated}.
\item Safety: \emph{executed-error rate} (share of sub-goals on which a wrong answer or action was
executed), executed precision, abstention or review rate, \emph{attack success rate} under injection,
and on AppWorld the \emph{collateral-damage rate}, the share of tasks failing a test that passes when the
agent does nothing.
\item Task success on AppWorld: the official task goal completion (TGC) and scenario goal completion (SGC).
\item Efficiency: tokens and cost, wall-clock latency, router mix and regret.
\item Context: share of tokens removed, fidelity $\Phi$, and the share of steps meeting the floor.
\item Certification: Learn-then-Test threshold, held-out risk against $\alpha$, coverage, and ECE.
\end{itemize}

\section{Statistical Analysis and Sample Size}
\label{sec:method-stats}

On ToolBench and KCC the unit of analysis is the session; each headline comparison is tested with a
paired bootstrap 95\,\% confidence interval \citep{efron1993bootstrap} and a Wilcoxon signed-rank test
\citep{wilcoxon1945individual}, and proportions use Wilson intervals \citep{wilson1927probable}. On
AppWorld the unit is the task: configurations run on the same backbone model are compared by McNemar's
exact test on per-task success \citep{mcnemar1947note} and by a paired bootstrap interval on the
difference in TGC; final configurations are run with at least three seeds and reported with pass$^k$.
One comparison per objective is declared primary; all others are exploratory.

\textbf{Is the sample large enough?} For O1 on ToolBench the per-session NDCG@5 difference between SATR
v2 and BM25 has $d_z = 0.436$; with 100 sessions the achieved power at $\alpha = 0.05$ is 0.99, and the
minimum detectable effect at 80\,\% power is 0.044 NDCG@5. For O3 the change in accepted errors from
24.0\,\% to 0.6\,\% has Cohen's $h = 0.87$ \citep{cohen1988power}. For certified execution, a threshold
can be certified at risk $\alpha$ and level $\delta$ only if $n \ge \lceil\log\delta/\log(1-\alpha)\rceil$
executions are available: 149 for $\alpha = 0.02$ and 59 for $\alpha = 0.05$ at $\delta = 0.05$. A
calibration half of ToolBench has about 87, which is why the stronger claims are made on the larger KCC
and AppWorld layers. On AppWorld, a pilot on \texttt{dev} estimates the discordant-pair rate between
configurations before test runs are spent, so that McNemar's test has adequate power.

\section{Threats to Validity}
\label{sec:method-threats}

\begin{itemize}[leftmargin=*]
\item \textbf{Simulated agents} in the ToolBench and KCC layers measure mechanisms, not LLM quality; the
AppWorld layer uses real open-weight agents.
\item \textbf{Pseudo-sessions and label noise} in KCC-Sessions; KCC services other than the
data.gov.in resources are simulated.
\item \textbf{Exploration on AppWorld}: outcome-level consensus executes candidate programs from a saved
world state and restores it; these executions are disclosed and counted, and a real deployment needs an
equivalent sandbox or an approval step.
\item \textbf{Model scale}: free and low-cost GPUs limit the backbone to at most about 32B parameters,
so absolute AppWorld scores are compared with entries of similar model size, and frontier entries are
cited as context.
\item \textbf{Hand-set parameters} for O2--O4, examined in sensitivity sweeps and, for the certification
threshold, replaced by Learn-then-Test.
\end{itemize}

All code, splits and aggregate results are public at \url{https://github.com/joyjeni/phd-agentic-ai-master}
and \url{https://github.com/joyjeni/acrs-appworld}; AppWorld content is never redistributed unencrypted,
as its licence requires.
